\section{Related Work}
\label{sec:related_work}

\subsection{Urban Crime Forecasting}
\label{subsec:related_crime_forecasting}

Crime forecasting seeks to anticipate future incident counts, hotspot status,
or event risk at operationally relevant locations and times. \rev{This variation
in target is accompanied by variation in scale and model design. Reviews of
spatial forecasting \citep{kounadi2020review} and multiscale methods
\citep{du2023review} frame the geographic and temporal choices, while a
machine- and deep-learning review \citep{mandalapu2023review} surveys the
algorithmic choices. Their different scopes help explain why reported
forecast scores are difficult to compare without a common target and
evaluation design.}

\rev{The models themselves show how the field has addressed those choices. Urban
metrics enter statistical learning as contextual predictors
\citep{alves2018crime}. Attentive recurrent networks model temporal structure
\citep{huang2018deepcrime}; a spatiotemporal deep-learning model
\citep{angbera2023model} and geographically informed risk models
\citep{deng2023risk} extend the representation of spatiotemporal dependence.
Multiscale hotspot mapping \citep{jing2024multiscale} and graph-based learning
of correlations among crime types \citep{wang2025mragnn} address different
forms of spatial and cross-category structure. More recent transformer
prediction \citep{butt2025start} and transfer with multi-source fusion
\citep{cui2026transfer} broaden the model and data choices; a survey of crime
event analysis documents the range of fusion strategies
\citep{hu2023fusion}. Across these families, performance is reported through
hotspot accuracy, prediction-accuracy indices, precision and recall, F1
scores, and count losses such as mean absolute error (MAE), root mean
squared error (RMSE), and mean squared error (MSE).}
Reported performance depends not only on the forecasting
algorithm but also on the target definition, spatial and temporal resolution,
validation period, comparison model, and distribution shift. Multidensity and
multiscale studies further show that changing the unit of analysis can alter
the forecasting problem itself \citep{cesario2024multidensity,jing2024multiscale},
while recent evaluation frameworks identify limited standardization and
cross-design comparability as persistent concerns \citep{amarante2025steval}.
Closely related urban-computing studies illustrate the same design dependence:
recent work models long-horizon traffic with contextual knowledge
\citep{huo2025traffic}, adapts shared-bicycle demand models to spatial and
temporal heterogeneity \citep{zhao2025bike}, and provides a unified toolbox
for constructing and comparing spatiotemporal prediction services
\citep{fang2025uctb}.
Thus, a lower test error establishes that one specification performs
better under a particular evaluation design; it does not establish how close
that specification is to an information-constrained prediction limit.
The unresolved issue is not a shortage of forecasting models, but the absence
of a model-independent way to assess whether an added information source
changes the attainable error benchmark for an integer-valued target.

\subsection{Mobility Data in Crime Analytics}
\label{subsec:related_mobility_crime}

Routine activity theory and ambient-population research imply that crime
opportunities depend on the changing presence of potential targets, offenders,
and guardians, not only on residential population
\citep{cohen1979routine,malleson2016ambient}. \rev{Song et al.\ relate daily
mobility flows to crime targets \citep{song2019legal}; Kadar and Pletikosa
use large-scale mobility data for long-term prediction
\citep{kadar2018mining}; and Kadar et al.\ use location-platform flows to
examine crime pattern theory \citep{kadar2020mobility}. Wu et al.\ combine
mobility flows with deep learning for short-term prediction
\citep{wu2022mobility}; Rummens et al.\ test mobile-phone denominators and
forecasts \citep{rummens2021mobile}; and Schertz et al.\ examine neighborhood
street activity \citep{schertz2021street}.} A recent systematic review documents
both the rapid growth and the heterogeneous uses of mobile-phone data in crime
research \citep{okmi2023mobile}. At finer scales, mobile-phone occupancy has
been related to outdoor robbery and theft \citep{valente2024routines}, and a
four-city forecasting study reports gains from combining mobility, crime
history, and sociodemographic features \citep{albors2025mobility}.

This evidence supports testing mobility as added information, but mobility
sources are not interchangeable. Digital traces differ in population
coverage, spatial support, temporal rhythm, directionality, and mode of travel
\citep{luca2023mobilitysurvey,zhao2024directionality}. Mobile-location data
can also exhibit scale- and time-dependent representation bias
\citep{li2024mobilitybias}. More broadly, a recent survey of foundation
models for the Internet of Things identifies context awareness, efficiency,
and safety as distinct criteria for incorporating sensor information
\citep{wei2026iot}. We therefore treat shared-bicycle records as a
high-frequency, spatially resolved, but mode-specific proxy for observed
public-space activity \citep{xu2023bike,ivanovic2023bike}, not as a census of
ambient population. Predictive information also does not establish a causal
effect: recent instrumental-variable evidence illustrates that predictive
association and causal identification can lead to different conclusions
\citep{vachuska2025causal}.

Most mobility--crime studies evaluate association patterns or changes in
fitted prediction scores. These questions remain useful, but they do not
separate the information value of a mobility source under a specified target
and information set from the ability of a particular model to exploit it.
Nor do they generally diagnose whether an apparent gain is induced by sparse states after mobility
variables are discretized and joined to a fine-scale crime panel. This gap
motivates the matched baseline--augmented information sets and finite-sample
calibration used in the present study.

\subsection{Information-Theoretic Prediction Limits}
\label{subsec:related_information_limits}

Conditional entropy measures uncertainty after admissible information has
been observed \citep{shannon1948,cover_thomas}. Entropy and Fano-type
arguments have been used to study potential predictability in human mobility
\citep{song2010limits}; conditional differential entropy has also been used to
bound traffic-forecasting performance \citep{li2022traffic}. Fang et al.\
derive general entropy-based variance bounds for continuous estimation and
prediction errors \citep{fang_variance_bounds}. Recent work has extended this
continuous line to time-series complexity ranking and entropy-based lower
bounds on predictive representation dimension
\citep{ayers2025complexity,hauser2026dmd}. These studies separate achieved
performance from an information-constrained limit, but their target spaces
and bound objects do not directly provide the integer-lattice MSE lower bound
required here.

Crime counts and many other operational targets are integer valued. A
continuous Gaussian argument cannot be made discrete by replacing
differential entropy with Shannon entropy: the error lies on a lattice, and
its maximum-entropy law under a second-moment constraint is a discrete
Gaussian. Integer maximum-entropy and lattice-correction results are
established ingredients \citep{massey_integer_entropy,rioul_massey}; general
Bayes-risk, rate--distortion, and Ziv--Zakai theories also show that lower
bounds on estimation risk are not new as a class. \rev{Shannon formulated
rate--distortion bounds \citep{shannon_rate_distortion}; Chen et al.\ studied
Bayes-risk bounds \citep{chen_bayes_risk}; and Jeong et al.\ examined
Ziv--Zakai bounds on MMSE \citep{jeong_ziv_zakai}.} What remains
underdeveloped for the present problem is an operational framework for
non-anticipating integer prediction that connects an explicit lattice MSE
floor to the predictor used, states necessary and sufficient attainability
conditions, diagnoses why the floor is not attained, values added
information, and calibrates its finite-sample estimate.

We refer to the MSE lower bound obtained by inverting the exact integer-lattice
maximum-entropy envelope as the \emph{Discrete Entropy Lower Bound} (DELB).
The theoretical package combines a universal floor, a predictor-dependent
refinement, strict attainability conditions, and a split of nonattainment
into residual history dependence and error-shape mismatch. To our knowledge,
these components have not previously been connected to incremental
information valuation and finite-sample calibration in a single operational
workflow for integer-valued urban forecasting. The contribution lies in this
integration and specialization, not in the discrete-Gaussian extremizer, the
lattice correction, or a new primitive divergence identity.

Table~\ref{tab:core_bound_comparison} makes the distinction operational by
comparing target, loss, support, and the role of each framework in the present
study.

\begin{table}[H]
\caption{Operational comparison of prediction-risk lower-bound frameworks.
The DELB contribution is an integrated specialization and calibration
workflow, not a new discrete-Gaussian maximum-entropy result.}
\label{tab:core_bound_comparison}
\centering
\scriptsize
\setlength{\tabcolsep}{2.5pt}
\begin{tabularx}{\textwidth}{
    >{\RaggedRight\arraybackslash}p{0.13\textwidth}
    >{\RaggedRight\arraybackslash}p{0.15\textwidth}
    >{\RaggedRight\arraybackslash}p{0.14\textwidth}
    >{\RaggedRight\arraybackslash}p{0.13\textwidth}
    >{\RaggedRight\arraybackslash}X}
\toprule
\textbf{Framework} & \textbf{Target} & \textbf{Loss or quantity} &
\textbf{Support} & \textbf{Relationship to this study}\\
\midrule
Fano inequality \citep{cover_thomas,feder_merhav}
& Finite classes & Hamming error & Finite label set
& Secondary classification-loss comparison; not an MSE bound\\

Fang et al.\ \citep{fang_variance_bounds}
& Continuous estimation or prediction & Variance/MSE & Continuous
& Continuous analogue; it does not directly handle an integer-lattice
error\\

Massey/Rioul \citep{massey_integer_entropy,rioul_massey}
& Integer-valued variable & Variance or second moment & Integer lattice
& Supplies the established maximum-entropy and lattice-correction basis\\

Bayes-risk/rate--distortion bounds
\citep{shannon_rate_distortion,chen_bayes_risk}
& General decision problem & General loss or distortion & General
& Broader than DELB, but often less explicit for the present count setting\\

Present DELB framework
& Non-anticipating integer prediction & MSE & Integer lattice
& Integrates an explicit floor, predictor refinement, attainability and
nonattainment diagnostics, added-information ordering, and finite-sample
calibration\\
\bottomrule
\end{tabularx}
\end{table}


Across these three literature streams, a specific applied methodological gap
remains. Crime forecasting provides increasingly flexible models,
and mobility studies provide increasingly detailed urban signals, whereas
existing entropy bounds do not supply a finite-sample-calibrated,
integer-lattice benchmark for determining whether such a signal changes the
information-constrained lower bound under a specified count target and information
set. DELB and the accompanying calibration
workflow address this intersection. The comparison is intentionally between
information sets rather than between algorithm families, and between
predictive evidence rather than causal effects.
